% Created 2023-01-27 Fri 16:15
% Intended LaTeX compiler: pdflatex
\documentclass[11pt]{article}
\usepackage[utf8]{inputenc}
\usepackage[T1]{fontenc}
\usepackage{graphicx}
\usepackage{longtable}
\usepackage{wrapfig}
\usepackage{rotating}
\usepackage[normalem]{ulem}
\usepackage{amsmath}
\usepackage{amssymb}
\usepackage{capt-of}
\usepackage{hyperref}
\author{mahmooz}
\date{\today}
\title{two dimensional random variable}
\hypersetup{
 pdfauthor={mahmooz},
 pdftitle={two dimensional random variable},
 pdfkeywords={},
 pdfsubject={},
 pdfcreator={Emacs 30.0.50 (Org mode 9.6)}, 
 pdflang={English}}
\begin{document}

\maketitle
\tableofcontents

\begin{definition}
the \href{discrete_maths2.org}{pair} \((X,Y)\) is a \href{20230104221353-random_variable.org}{random variable} if \(X\) and \(Y\) are defined over the same \href{probability.org}{probability space}\\[0pt]
refer to \href{20230120140031-joint_probability_distribution.org}{joint probability distribution}\\[0pt]
\begin{characteristic}
let \((X,Y)\) be a two dimensional random variable, the probability function (called \href{20230123011904-marginal_distribution.org}{marginal distribution}) of \(X\) is defined as \[ P(X=x)=\sum_{y} P(X=x,Y=y) \] and of \(Y\) is defined as \[P(Y=y)=\sum_{x} P(X=x,Y=y)\]\\[0pt]
\end{characteristic}
\begin{characteristic}
let \((X,Y)\) be a two dimensional random variable, the conditional probability function of \(X\) for \(Y=y\) is \[ P(X=x\mid Y=y)=\frac{P(X=x,Y=y)}{P(Y=y)} \] and similarly, \[ P(Y=y\mid X=x)=\frac{P(X=x,Y=y)}{P(X=x)} \]\\[0pt]
\end{characteristic}
\begin{my_example}
we throw a symmetric coin with the values 0 and 1, if we get 0 we throw a die, if we get 1 we throw the same coin again.\\[0pt]
we define \(X\) as the result of the first throw and \(Y\) as the result of the second\\[0pt]
we build a \href{20230123141739-joint_probability_table.org}{joint distribution table}\\[0pt]
\begin{center}
\begin{tabular}{rrlrrrrrl}
x\y & 0 & 1 & 2 & 3 & 4 & 5 & 6 & P(X=x)\\[0pt]
\hline
0 & 0 & 1/12 & 1/12 & 1/12 & 1/12 & 1/12 & 1/12 & 1/2\\[0pt]
1 & 1/4 & 1/4 & 0 & 0 & 0 & 0 & 0 & 1/2\\[0pt]
P(Y=y) & 1/4 & 4/12 & 1/12 & 1/12 & 1/12 & 1/12 & 1/12 & 1\\[0pt]
\end{tabular}
\end{center}
what is the probability that in the first throw we got 0 if we know in the second throw we got a number smaller than 3?\\[0pt]
\begin{answer}
\begin{gather*}
P(X=0|Y<3)=\frac{0+\frac{1}{12}+\frac{1}{12}}{\frac{1}{4}+\frac{4}{12}+\frac{1}{12}}\\
E(X)=0 \cdot \frac{1}{2} + 1 \cdot \frac{1}{2} = \frac{1}{2}\\
E(Y)=0 \cdot \frac{1}{4} + 1 \cdot \frac{4}{12} + 2 \cdot \frac{1}{12} + 3 \cdot \frac{1}{12} + 4 \cdot \frac{1}{12} + 5 \cdot \frac{1}{12} + 6 \cdot \frac{1}{12} = 2\\
\mathrm{Var}(X)=E(X^2)-(E(X))^2=0^2 \cdot \frac{1}{2} + 1^2 \cdot \frac{1}{2} - \left(\frac{1}{2}\right) = \frac{1}{4}\\
\vdots
\end{gather*}
\end{answer}
\end{my_example}
\begin{characteristic}
given \((X,Y)\) is a two dimensional random variable, \(X\) and \(Y\) are independent if\\[0pt]
\[
P(X=x,Y=y)=P(X=x) \cdot P(Y=y)
\]\\[0pt]
for all values of \(X\) and \(Y\)\\[0pt]
\end{characteristic}
\begin{characteristic}
let \((X,Y)\) be a two dimensional random variable, let \(f(X,Y)\) be a function of \(X\) and \(Y\), then\\[0pt]
\[
E(f(X,Y)) = \sum_{x} \sum_{y} f(x,y) \cdot P(X=x,Y=y)
\]\\[0pt]
\end{characteristic}
\begin{characteristic}
let \((X,Y)\) be a two dimensional random variable\\[0pt]
we know \(E(X+Y)=E(X)+E(Y)\), we denote \(E(Y)=\mu_2,E(X)=\mu_1\)\\[0pt]
\begin{align*}
  \mathrm{Var}(X+Y) &= E\left(\left(X+Y-\mu_1-\mu_2\right)^2\right)\\
  &= E\left((X-\mu_1)^2+(Y-\mu_2^2)+2(X-\mu_1)(Y-\mu_2)\right)\\
  &= E\left((X-\mu_1)^2\right)+E\left(\left(Y-\mu_2\right)^2\right)+2E((X-\mu_1)(Y-\mu_2))\\
  &= \mathrm{Var}(X)+\mathrm{Var}(Y)+2E((X-E(X))(Y-E(Y)))
\end{align*}
\end{characteristic}
\end{definition}
\end{document}